\documentclass[11pt]{article}
\usepackage{amsmath,amssymb,bm,mathtools}
\usepackage[margin=0.9in,a4paper]{geometry}
\usepackage{booktabs,array,multirow}
\usepackage{xcolor}
\usepackage[most]{tcolorbox}
\usepackage[colorlinks=true,linkcolor=blue!70!black,urlcolor=blue!70!black]{hyperref}
\usepackage{titlesec}
\usepackage{listings}
\usepackage{microtype}

\definecolor{myblue}{RGB}{0,0,180}
\definecolor{lightblue}{RGB}{214,234,255}
\definecolor{darkblue}{RGB}{10,40,120}
\definecolor{darkorange}{RGB}{180,90,0}
\definecolor{lightorange}{RGB}{255,237,200}
\definecolor{lightgreen}{RGB}{222,245,222}
\definecolor{darkgreen}{RGB}{20,100,40}
\definecolor{hl}{RGB}{200,30,30}

\titleformat{\section}{\large\bfseries\color{myblue}}{\thesection}{0.8em}{}
\titleformat{\subsection}{\normalsize\bfseries\color{myblue}}{\thesubsection}{0.8em}{}
\titleformat{\subsubsection}{\normalsize\bfseries\color{myblue}}{\thesubsubsection}{0.8em}{}

\newtcolorbox{conventionbox}{
  colback=lightblue, colframe=darkblue,
  title={\textbf{Convention --- read this first}},
  fonttitle=\bfseries, coltitle=white, colbacktitle=darkblue,
  boxrule=0.8pt, arc=2pt}
\newtcolorbox{notebox}[1]{
  colback=lightblue, colframe=darkblue,
  title={\textbf{#1}},
  fonttitle=\bfseries, coltitle=white, colbacktitle=darkblue,
  boxrule=0.8pt, arc=2pt}
\newtcolorbox{warningbox}[1]{
  colback=lightorange, colframe=darkorange,
  title={\textbf{#1}},
  fonttitle=\bfseries, coltitle=white, colbacktitle=darkorange,
  boxrule=0.8pt, arc=2pt}
\newtcolorbox{interpbox}[1]{
  colback=lightgreen, colframe=darkgreen,
  title={\textbf{#1}},
  fonttitle=\bfseries, coltitle=white, colbacktitle=darkgreen,
  boxrule=0.8pt, arc=2pt}

\lstset{language=Python, basicstyle=\small\ttfamily,
  commentstyle=\color{green!60!black}, keywordstyle=\color{blue},
  showstringspaces=false, columns=flexible, breaklines=true}

\newcommand{\R}{\mathbb{R}}
\newcommand{\Wqk}{W_{QK}}
\newcommand{\Wov}{W_{OV}}
\newcommand{\Wq}{W_Q}
\newcommand{\Wk}{W_K}
\newcommand{\Wv}{W_V}
\newcommand{\Wo}{W_O}
\newcommand{\sd}{\sqrt{D}}
\newcommand{\rowlab}[1]{\ \leftarrow\ \text{\footnotesize #1}}
\newcommand{\collab}[1]{\text{\footnotesize #1}}

\newcommand{\Vsig}{V_{\mathrm{sig}}}
\newcommand{\nnoise}{n_{\mathrm{noise}}}
\newcommand{\bsig}{b_{\mathrm{sig}}}
\newcommand{\bnoise}{b_{\mathrm{noise}}}

\title{\textbf{Hand construction, reduced-magnitude form}\\[2pt]
\large Combined $\Wqk$ / $\Wov$, no unembedding, no log, no output softmax}
\author{ICL --- Markov-with-noise}
\date{\today}

\begin{document}
\maketitle

\begin{conventionbox}
$V=8$ tokens: \textbf{signal} $0..3$ ($\Vsig=4$), \textbf{noise} $4..7$ ($\nnoise=4$).
$T=100$ positions, scored on $[50,100)$. Residual widths
$D_1=V+T=108$ (layer-1 input), $D_2=2D_1=216$ (layer-2 input).
Attention scores are $s_{ij}=x_i^\top \Wqk\, x_j/\sqrt{D_1}$, so \textbf{a $\Wqk$ entry of size $w$
contributes $w/\sqrt{108}\approx w/10.39$ to the logit}. Everything below is stated first in
\emph{logit} units and then multiplied by $\sqrt{D_1}$ to give the matrix entry.
\end{conventionbox}

\section{What changed, and why}

This document supersedes the earlier combined construction. The circuit is identical; three
things about the \emph{numbers} changed.

\begin{center}\small
\begin{tabular}{@{}lrrl@{}}
\toprule
& old & new & why \\
\midrule
$|\Wqk^{(1)}|_{\max}$ & $1.04\times10^{10}$ & $1039.2$ & see below \\
$|\Wqk^{(2)}|_{\max}$ & $5.20\times10^{6}$  & $1039.2$ & see below \\
nonzeros in $\Wqk^{(2)}$ & 192 & \textbf{40} & inert entries removed \\
\bottomrule
\end{tabular}
\end{center}

\subsection{Change 1 --- magnitudes reduced to the smallest that still works}

The old constants ($P_1=1000$, $\Gamma=-5\times10^{4}$, $\text{BIG}=-10^{9}$) were chosen to make
every softmax exactly one-hot with enormous margin. They are far past what is needed: a logit gap
of $30$ already gives $e^{-30}\approx10^{-13}$ leakage. They also put the construction in a
magnitude regime gradient descent never visits, which makes any entry-by-entry comparison with a
trained model meaningless.

The reduction was measured, not guessed. Sweeping the logit-level constants and re-evaluating the
cross entropy gives

\begin{center}\small
\begin{tabular}{@{}rrrrrr@{}}
\toprule
$\lambda$ & $K$ & $B$ & $|\Wqk^{(1)}|_{\max}$ & $|\Wqk^{(2)}|_{\max}$ & CE $\alpha{=}0.6$ \\
\midrule
20 & 1000 & $10^9$ & $1.0\times10^{10}$ & $5.2\times10^{6}$ & 1.69622 \ \small(original) \\
\textbf{20} & \textbf{100} & \textbf{100} & \textbf{1039.2} & \textbf{1039.2} & \textbf{1.70035} \\
20 & 100 & 30  & 1039.2 & 1039.2 & 1.87644 \\
20 & 100 & 15  & 1039.2 & 1039.2 & 3.00135 \\
20 & 30  & 100 & 1039.2 & 1039.2 & 1.73458 \\
6  & 12  & 20  & \ \ 311.8 & \ \ 207.8 & 1.77789 \\
\bottomrule
\end{tabular}
\end{center}

So $(\lambda,K,B)=(20,100,100)$ costs \textbf{0.004 nats} against the original while shrinking
$\Wqk^{(2)}$ by a factor of $5000$ and $\Wqk^{(1)}$ by $10^{7}$.

\begin{warningbox}{The veto depth $B$ is the binding constraint --- not the induction bonus $K$}
Dropping $K$ from 100 to 30 costs 0.034 nats. Dropping $B$ from 100 to 30 costs 0.176, and $B=15$
is catastrophic (3.00). The reason: a noise key that happens to satisfy the induction match
receives $+K$ from the induction block and $-B$ from noise suppression. If $B \not> K$ the veto
loses and noise leaks into the count. \textbf{Keep $B \ge K$.}
\end{warningbox}

\begin{notebox}{How this compares with what gradient descent finds}
The trained run (Adam, lr 1.0, zero init, 500 steps, $\alpha=0.6$) reaches CE $1.6965$ with
$|\Wqk^{(1)}|_{\max}=110.9$ and $|\Wqk^{(2)}|_{\max}=67.4$ --- roughly $10\times$ smaller still.
Its induction diagonal is $+67$ against a background of $-55$, i.e. about $12$ logits of
separation, where this construction uses $100$. So GD reaches the same loss with a much softer
circuit. Pushing the construction to GD's magnitudes ($\lambda,K,B$ near $6,12,20$) costs
$0.08$ nats, so the two are \emph{not} the same solution at different temperatures --- GD is doing
something genuinely more economical, and that gap is a result worth reporting rather than tuning
away.
\end{notebox}

\subsection{Change 2 --- the inert column in the induction block is removed}

The old $\Wqk^{(2)}[0{:}8,\,108{:}116]$ carried $+103.92$ in \emph{every} row of column 4
(the first noise token). It was never a design choice: it is a factorisation side effect. The
single entry $\Wq^{(2)}[4,i]=\sqrt{D_1}\cdot\text{NF}$ appears twice in $\Wqk=\Wq^\top\Wk$,
\[
\underbrace{\times\, \Wk^{(2)}[4,\,D_1{+}4]=1}_{\textstyle +103.92 \text{ in the induction block}}
\qquad\qquad
\underbrace{\times\, \Wk^{(2)}[4,\,4{:}8]=\Gamma}_{\textstyle -5.2\times10^{6}\text{ noise suppression}}
\]
so the noise-suppression term could not be written without also emitting that column. It does
nothing, for two reasons: it is $100\times$ below the induction diagonal, and the keys it boosts
have already been vetoed at $-5.2\times10^{6}$. \textbf{Ablating it changes the cross entropy by
exactly $0.00\times10^{0}$ at every $\alpha$} --- bit-identical, while ablating noise suppression
for comparison costs $0.133$ nats. It is now zero.

\subsection{Change 3 --- the masked-out future entries are zeroed}

The positional block was filled with $-\lambda(t_q-t_k)$ for \emph{all} pairs. For $t_k>t_q$ that
is positive and grows to $+\lambda\cdot99$, which is where $|\Wqk^{(1)}|_{\max}=20577$ came from
even after capping. Those entries sit strictly above the causal mask and are discarded by
$\text{softmax}(s+M)$ every single forward pass. Setting them to zero costs nothing measurable and
removes the largest numbers in the matrix. The past decay is additionally \textbf{capped at
$\Delta=5$}: at $\alpha=0.6$ the chance of no signal token in five steps is $0.4^5\approx1\%$, so
ranking beyond that is wasted range.

\section{The construction, matrix by matrix}

Logit-level constants: $\lambda=20$ (positional slope), $C=5$ (decay cap), $K=100$ (induction
bonus), $B=100$ (veto depth). Matrix entries are these times $\sqrt{D_1}=10.3923$.

\subsection{Layer 1 --- previous-signal retriever, $\Wqk^{(1)}\in\R^{108\times108}$}

Job: from position $i$, attend to the \textbf{most recent signal token strictly before $i$}.
That token is the state the next emission comes from.

\[
\Wqk^{(1)}[r,c]=\sqrt{D_1}\cdot\begin{cases}
-B & r\in[4,8),\ \text{any } c \rowlab{noise \emph{query} rows: vetoed}\\[2pt]
-B & r=V{+}t,\ c=V{+}t \rowlab{exclude self, $t_k=t_q$}\\[2pt]
-\lambda\min(t_q-t_k,\,C) & r=V{+}t_k,\ c=V{+}t_q,\ t_k<t_q \rowlab{capped recency}\\[2pt]
0 & t_k>t_q \rowlab{above the causal mask}\\[2pt]
0 & \text{elsewhere}
\end{cases}
\]

\begin{center}\small
\begin{tabular}{@{}llrr@{}}
\toprule
block & role & nonzeros & $|\cdot|_{\max}$ \\
\midrule
$[0{:}8,\ 0{:}8]$      & token $\times$ token, noise-query veto & 32 & 1039.2 \\
$[8{:}108,\ 0{:}8]$    & position $\times$ token & 400 & 1039.2 \\
$[8{:}108,\ 8{:}108]$  & position $\times$ position, capped recency + self-veto & 5050 & 1039.2 \\
$[0{:}8,\ 8{:}108]$    & (empty) & 0 & --- \\
\midrule
\textbf{total} & & \textbf{5482} of 11664 & \textbf{1039.2} \\
\bottomrule
\end{tabular}
\end{center}

\subsection{Layer 1 --- $\Wov^{(1)}\in\R^{108\times108}$}
$\Wov^{(1)}[0{:}8,\,0{:}8]=I_8$, zero elsewhere. Attention copies the \emph{token one-hot} of
whichever position it selected, so layer 1 writes a one-hot of the retrieved state into the first
8 coordinates of its output. Frozen, never trained.

\subsection{Layer 2 --- counting induction head, $\Wqk^{(2)}\in\R^{216\times216}$}

Layer 2's input is $x^{(1)}=[\,x^{(0)}\mid z^{(1)}\,]$, so columns $0{:}8$ are the raw token
one-hot, $8{:}108$ the position one-hot, and $108{:}116$ the state layer 1 retrieved.

\[
\Wqk^{(2)}[r,c]=\sqrt{D_1}\cdot\begin{cases}
+K & r=v,\ c=D_1{+}v,\ v<\Vsig \rowlab{\textbf{induction}: own token $=$ key's state}\\[2pt]
-B & r<V,\ c\in[4,8) \rowlab{\textbf{noise suppression}: veto noise keys}\\[2pt]
+K/20 & r=D_1{+}v,\ c=D_1{+}v,\ v<\Vsig \rowlab{own state $=$ key's state}\\[2pt]
0 & \text{everywhere else}
\end{cases}
\]

\begin{center}\small
\begin{tabular}{@{}llrr@{}}
\toprule
block & role & nonzeros & $|\cdot|_{\max}$ \\
\midrule
$[0{:}8,\ 0{:}8]$          & noise suppression & 32 & 1039.2 \\
$[0{:}8,\ 108{:}116]$      & induction head (diagonal only now) & \textbf{4} & 1039.2 \\
$[108{:}116,\ 108{:}116]$  & state $\times$ state & 4 & \ \ 52.0 \\
\midrule
\textbf{total} & & \textbf{40} of 46656 & \textbf{1039.2} \\
\bottomrule
\end{tabular}
\end{center}

The induction block is now exactly $\sqrt{D_1}K\,I_4$ padded with zeros --- \textbf{4 nonzeros
instead of 20}, because the artefact column is gone and the noise-query rows were always zero:
\[
\Wqk^{(2)}[0{:}8,\,108{:}116]=
\begin{pmatrix}
1039.2 & 0 & 0 & 0 & 0&0&0&0\\
0 & 1039.2 & 0 & 0 & 0&0&0&0\\
0 & 0 & 1039.2 & 0 & 0&0&0&0\\
0 & 0 & 0 & 1039.2 & 0&0&0&0\\
0&0&0&0&0&0&0&0\\ 0&0&0&0&0&0&0&0\\ 0&0&0&0&0&0&0&0\\ 0&0&0&0&0&0&0&0
\end{pmatrix}
\]

\subsection{Layer 2 --- $\Wov^{(2)}\in\R^{216\times108}$}
$\Wov^{(2)}[0{:}8,\,0{:}8]=I_8$, zero elsewhere. Frozen.

\begin{interpbox}{Why this one choice removes the unembedding, the log and the output softmax}
With $\Wov^{(2)}[0{:}8,\,0{:}8]=I$, the first $V$ coordinates of layer 2's attention output are
\[
f_v=\sum_j A^{(2)}_{ij}\,\mathbb{1}[x_j=v].
\]
The attention weights $A^{(2)}_{ij}$ are non-negative and sum to 1, and each $\text{tok}_j$ is a
one-hot. A convex combination of one-hots \textbf{is} a distribution: $f\ge0$ and $\sum_v f_v=1$
exactly. So $f$ is already a probability vector over the 8 tokens --- there is nothing for an
unembedding, a log, or a final softmax to do. That is the whole reason they are gone.

This guarantee is exactly as strong as ``$\Wov$ frozen to the identity''. If $\Wov$ is trained,
$f$ becomes a convex combination of \emph{arbitrary} vectors, $f\ge0$ fails, and everything below
breaks.
\end{interpbox}

\section{The readout bias --- why it exists and where its formula comes from}

This is the part of the design that is easiest to get wrong, so it is worked out slowly.

\subsection{Step 1 --- what layer 2 actually hands us}

With the noise filter working and the induction head perfect, layer 2's output at $\alpha=0.6$ is
\[
f=[\,\underbrace{0.5,\ 0.5,\ 0,\ 0}_{\text{signal slots }0..3}\ \mid\
\underbrace{0,\ 0,\ 0,\ 0}_{\text{noise slots }4..7}\,]
\]
The signal slots hold the successor distribution read off the context and sum to 1. The noise
slots are \textbf{exactly zero} --- that is what ``the noise filter works'' means: all of layer 2's
attention went to signal keys, so no noise one-hot got averaged in.

\subsection{Step 2 --- why that is a disaster without a bias}

The target token is a noise token $(1-\alpha)=40\%$ of the time. When it is, we need
$-\log p_{\text{noise}}$, and $p_{\text{noise}}=0$.

\begin{warningbox}{The trap}
Numerically the code floors it, so the cost is $-\log(\text{LOG\_FLOOR})=69$ nats on 40\% of
positions. Gradient descent will do anything to avoid that --- and the only thing available is to
\textbf{break the noise filter}: park $(1-\alpha)$ of layer 2's attention on noise positions so the
noise slots are non-zero. Measured, this is exactly what happens: with \texttt{readout=none} the
trained model sits at $\text{signal\_mass}=0.598\approx\alpha$ instead of $1.0$, and the loss
plateaus at \textbf{1.847} instead of 1.696.

So the missing bias does not just cost accuracy --- \emph{it forbids the circuit we are trying to
study.} One attention distribution is forced to do two jobs: estimate the successor counts
\emph{and} carry the noise mass.
\end{warningbox}

The fix is to supply the noise mass as a \textbf{known constant}. We know $\alpha$; the noise
distribution is uniform and context-independent. There is nothing to learn there. Handing it over
as a constant frees the entire attention budget for the one job that does need the context.

\subsection{Step 3 --- what we want the output to be}

At $\alpha=0.6$ the target is
\[
p=[\,\underbrace{0.3,\ 0.3,\ 0,\ 0}_{\alpha=0.6 \text{ split over the successors}}\ \mid\
\underbrace{0.1,\ 0.1,\ 0.1,\ 0.1}_{(1-\alpha)=0.4 \text{ spread over 4 noise tokens}}\,]
\]
The noise mass is missing from $f$ entirely, and $\bnoise$ has to supply it.

\subsection{Step 4 --- the one step that catches people out}

The readout is $p=(f+b)/\!\sum(f+b)$. \textbf{We divide by the sum}, so whatever we put in gets
shrunk. Work out the three sums:
\[
\begin{aligned}
\text{signal side} &= \textstyle\sum f + \Vsig\,\bsig = 1+4(0.01)=1.04\\
\text{noise side}  &= \nnoise\, b = 4b\\
\text{total}       &= 1.04+4b
\end{aligned}
\]
Now demand that the noise side is exactly $1-\alpha$ \emph{of the total}:
\[
\frac{4b}{1.04+4b}=0.4
\]
\[
\begin{aligned}
4b &= 0.4(1.04)+0.4(4b) &&\text{multiply out}\\
4b-1.6b &= 0.416 &&\text{collect } b \text{ on the left}\\
2.4\,b &= 0.416 &&\text{and } 4-1.6=4(0.6)=\nnoise\cdot\alpha\\
b &= 0.416/2.4 = 0.173333
\end{aligned}
\]

\begin{interpbox}{The step that matters: $4-1.6=2.4$}
That is $\nnoise\big(1-(1-\alpha)\big)=\nnoise\cdot\alpha$. \textbf{Collecting the $b$ terms is
what puts $\alpha$ in the denominator} --- it is not the $\alpha$ you started with. This is the
whole reason the answer is not the naive $(1-\alpha)/\nnoise$.
\end{interpbox}

Replacing the numbers by symbols --- $0.4\to(1-\alpha)$, $1.04\to S=1+\Vsig\bsig$,
$4\to\nnoise$, and the $4-1.6$ step $\to\nnoise\alpha$:
\[
\boxed{\ \bnoise=\frac{(1-\alpha)\,\bigl(1+\Vsig\,\bsig\bigr)}{\nnoise\cdot\alpha}\ }
\]

\subsection{Step 5 --- check it}

\[
(f+b)/\!\sum(f+b)=[\,0.2942,\ 0.2942,\ 0.0058,\ 0.0058\ \mid\ 0.1,\ 0.1,\ 0.1,\ 0.1\,]
\]
Noise slots: $0.100000$. Target $(1-\alpha)/\nnoise=0.100000$. Exact.

\subsection{Step 6 --- why not just $b=(1-\alpha)/\nnoise=0.1$?}

\[
(f+b)/\!\sum(f+b)=[\,0.3542,\ 0.3542,\ 0.0069,\ 0.0069\ \mid\ 0.0694\times4\,]
\]
The noise slots come out $0.0694$, not $0.1$; total noise mass $0.278$, not $0.4$.

\begin{interpbox}{One sentence to remember}
\textbf{Dividing by the sum shrinks whatever you put in, so you have to put in more.} Dividing by
$\alpha$ \emph{is} that correction: $0.1/0.6=0.1667$, and the remaining $\times1.04$ pays for the
$\Vsig\bsig$ we also added to the signal side. $0.1667\times1.04=0.1733$.
\end{interpbox}

\subsection{And $\bsig=0.01$ --- chosen, not derived}

$\bsig$ is a smoothing constant. It stops $f_v=0$ for a signal token the context never happened to
show as a successor, which would otherwise be $\log 0$. Classic add-$\beta$ smoothing with
$\beta=0.01$.

It is not free: in the check above the two unseen signal slots take $0.0058$ each, and the real
successors give up exactly that ($0.2942$ instead of $0.3000$). Note it appears \emph{inside} the
$\bnoise$ formula precisely so the noise side compensates for the mass it adds.

\begin{notebox}{$\alpha=1$}
$(1-\alpha)/(\nnoise\alpha)=0$, so $\bnoise$ would be zero and any noise target would cost
$\log 0$. There are no noise targets at $\alpha=1$, but to keep the slot off exactly zero the code
uses $\bnoise=10^{-9}$.
\end{notebox}

\section{After the bias: the average, then the loss}

Two averages happen, in this order, and they are different things.

\subsection{(a) The renormalisation --- the ``average'' inside the readout}
\[
p=\frac{f+b}{\sum_v (f+b)_v}
\]
Dividing an 8-vector by its own sum makes it a distribution again. \textbf{This is the step that
replaced the old $\log\to\text{softmax}$ pair}, because
$\mathrm{softmax}(\log x)=x/\!\sum x$ exactly --- they were always the same operation. Measured
difference between the two: $5.6\times10^{-17}$.

\subsection{(b) The mean over scored tokens --- the average before the reported loss}
\[
\mathcal{L}=-\frac{1}{|B|\cdot(T_{\text{end}}-T_{\text{start}})}
\sum_{s\in B}\ \sum_{i=T_{\text{start}}}^{T_{\text{end}}-1}\log p^{(s)}_{i,\,y^{(s)}_i}
\]
with $T_{\text{start}}=50$, $T_{\text{end}}=100$: the first 50 positions are context only and are
never scored. In code this is the single \texttt{.mean()} at the end of \texttt{ce\_loss}, taken
over sequences \emph{and} positions together. Training is full batch, so the gradient is
accumulated over chunks with each chunk's mean reweighted by its share of the scored tokens ---
arithmetically identical to one backward pass over all 25{,}000 sequences.

\begin{notebox}{Why the loss is now well behaved}
$f\ge0$ (convex combination of one-hots) and $b>0$ give $p$ strictly inside $(0,1)$. So
$\mathcal{L}\ge0$ and the numerical floor never binds --- measured margin $\sim10^{28}$. Under
\texttt{readout=none} the floor is load-bearing: it binds on $0.098\%$ of positions at
$\alpha=0.6$. \textbf{Both guarantees require $\Wov$ frozen to the identity.} With $\Wov$
trainable, $f$ can go negative, $\sum(f+b)$ can pass through zero, and $-\log p$ becomes unbounded
below --- this produced $\text{best\_train}=-1.05$ before the guard existed.
\end{notebox}

\section{Summary of every matrix}

\begin{center}\small
\begin{tabular}{@{}llrrl@{}}
\toprule
matrix & shape & nonzeros & $|\cdot|_{\max}$ & trained? \\
\midrule
$\Wqk^{(1)}$ & $108\times108$ & 5482 & 1039.2 & yes (2 blocks, 10064 free) \\
$\Wov^{(1)}$ & $108\times108$ & 8 & 1.0 & no --- frozen $I_8$ \\
$\Wqk^{(2)}$ & $216\times216$ & \textbf{40} & 1039.2 & yes (3 blocks, 192 free) \\
$\Wov^{(2)}$ & $216\times108$ & 8 & 1.0 & no --- frozen $I_8$ \\
\midrule
$b$ & $8$ & 8 & 0.1733 & no --- computed from $\alpha$ \\
\bottomrule
\end{tabular}
\end{center}

There is no unembedding, no log, and no output softmax. Cross entropy at $\alpha=0.6$:
\textbf{1.70035} for this construction, $1.69622$ for the original large-magnitude one,
$1.6965$ for the trained model, and $1.6434$ for ideal-$\alpha$ (which knows the true transition
table and is therefore unreachable from context alone).

\section{Next: the ICL score}

This construction is the reference for the \emph{ICL score}: the cosine between a trained
$\Wqk$ and this one, restricted to the entries gradient descent is allowed to move,
\[
\text{ICL score}_\ell=\cos\Bigl(\Wqk^{(\ell)}\big|_{\text{mask}},\ \
\Wqk^{(\ell),\text{constr}}\big|_{\text{mask}}\Bigr).
\]
Two reasons the reduced construction is the right reference. First, the old one put
$\sim99.98\%$ of its layer-2 energy in 4 diagonal entries of size $10^4$ and the rest in an inert
artefact column, so the cosine was dominated by numbers that do no work. With 40 nonzeros and one
magnitude, every entry that contributes is an entry that matters. Second, the old magnitudes were
$10^7$ away from anything GD produces; a comparison against them measures the gauge, not the
circuit.

\begin{warningbox}{The ceiling to quote alongside any ICL score}
The trainable mask holds \textbf{100\%} of $\Wqk^{(2)}$'s construction energy --- its 192 free
cells are exactly the construction's layer-2 support --- but only \textbf{24.8\%} of
$\Wqk^{(1)}$'s, because the construction puts mass on position-rows $\times$ noise-token-columns
which the block-sparse mask excludes. So layer 2's score is comparable against $1.0$; layer 1's is
capped far below it, and a small number there is not a failure.
\end{warningbox}

\end{document}
